\documentstyle[% 


righttag% 


,11pt% 


,amssymb% 


,russian%               remove in Eng. 


,izvru%                 Set izven% in Eng. 


% ,izvru-ks             for brief communications 


]{amsart} 


 


%       Math definitions 


 


\newcommand{\tts}[1]{}


\newcommand{\al}{\alpha}


\newcommand{\be}{\beta}


\newcommand{\sig}{\sigma}


\newcommand{\fo}{\forall}


\newcommand{\pli}{+\infty}





\newcommand{\En}[1]{\bold E_{#1}}


\newcommand{\Arg}{\operatorname{Arg}}


\newcommand{\opt}{\operatorname{opt}}


\newcommand{\R}{\bold R}


\newcommand{\F}{\bold F}


\newcommand{\X}{\bold X}


\newcommand{\sij}[2]{\sum_{#1}^{#2}}


\newcommand{\svij}[2]{\sum\limits_{#1}^{#2}}


\newcommand{\ij}[1]{1,\dotsc,#1}


\newcommand{\nk}[2]{#1,\dotsc,#2}





\newcommand{\bi}{\bigtriangleup}


\newcommand{\Lra}{\Longrightarrow}


\newcommand{\Llra}{\Longleftrightarrow}


\newcommand{\su}{\subset}


\newcommand{\emp}{\emptyset}





\theoremstyle{plain} 


\theorembodyfont{\it} 


\newtheorem{thms}{\hskip\parindent Теорема}[section]


\newtheorem{lems}{\hskip\parindent Лемма}[section]


 


\theoremstyle{definition} 


\theorembodyfont{\rm} 


\newtheorem{cors}{\hskip\parindent Следствие}[section]


\newtheorem{notes}{\hskip\parindent Замечание}[section]


\numberwithin{equation}{section}


 


%\hoffset=-15mm%                  For Eng. 


 


\setcounter{page}{49} 


\god{1997} 


\nomer{12~(427)} %                Remove in Eng. 


%\nomer{Vol.\,41, No.\,12}%        Set in Eng. 


% \str{75--81}%                    Set in Eng. 


\udk{519.853} 


\author{\it И.И.\,Еремин} 


\title{Некоторые вопросы\\


кусочно-линейного программирования} 


 


\thanks{Работа выполнена при


финансовой поддержке Российского фонда фундаментальных исследований


(проект 96--01--00116).} 


\date{} 


% \pagestyle{myheadings}%         Set in Eng. 


 


\pagestyle{plain} %              Remove in Eng. 


\begin{document} 


 


% \thispagestyle{plain}%          Set in Eng. 


 


\maketitle 


 


Задачи кусочно-линейной оптимизации составляют важный класс


задач математического программирования. С одной стороны, они


ближайшим образом обобщают задачи линейного программирования,


сохраняя многие свойства последних, а с другой --- служат


самостоятельным хорошо работающим аппаратом в различных областях


прикладной математики. Кусочно-линейными задачами можно


аппроксимировать нелинейные задачи оптимизации, бывает полезным


редуцировать к ним те или иные типы многоэкстремальных задач. Есть


еще одно важное обстоятельство, характеризующее кусочно-линейные


функции: они обладают уникальными разделительными свойствами, что


позволяет их использовать в дискриминантном анализе.





Алгебра кусочно-линейных функций ($k$-{\it функций}) и


задач кусочно-линейного программирования может изучаться


в рамках более общих конструкций, а именно, --- в


рамках $\sig$-расширений функциональных пространств. Речь идет об


алгебраическом расширении фиксированного функционального


пространства $\F_0$, замкнутого относительно операции {\it


дискретного максимума}. Если $\F$ --- указанное расширение, а


$\F_0$ --- пространство линейных функций, то $\F$ --- пространство


кусочно-линейных функций; если $\F_0$ --- класс квадратичных


функций, то $\F$ --- класс кусочно-квадратичных функций и т.д.





Хотя кусочно-линейные функции и соответствующий аппарат для них


имеют важное значение, тем не менее число работ, посвященных этим


вопросам, незначительно. Отметим работы \cite{Plo}--\cite{Gor}, первая из


которых содержит


главу III под названием ``Кусочно-линейные функции и полиэдральные


множества''. В этой главе проводится довольно основательное


исследование по алгебре и геометрии класса кусочно-линейных функций.





\section{$\boldsymbol \sig$-расширения функциональных


пространств}





Пусть $\F_0$ --- некоторое функциональное 


пространство с вещественным пространством $\X$ значений 


аргумента функций из $\F_0$. Если $\{f_j(x)\}_{j\in


J} \su \F_0$, то функция {\it дискретного максимума}


$f(x):= \max\limits_{j\in J} f_j(x)$ может как


принадлежать $\F_0$, так и не принадлежать. Способ образования


функции $f(x)$ назовем $\sig$-{\it операцией}.





Речь пойдет о минимальном расширении пространства $\F_0$ до


пространства $\F$, обеспечивающего свойство $\sig$-{\it замкнутости}


\begin{equation}


\{f_j(x)\}_{j\in J} \su \F \,\Lra\, \max_{j\in J}


f_j(x) \in \F.


\end{equation}


В этой ситуации выполняется, очевидно, свойство линейной


замкнутости


\begin{equation}


\{f_j^i \mid i\in I,\; j\in J_i \} \su \F \,\Lra\,


\sij{i\in I}{} \al_i\, \max_{j\in J_i}\, f_j^i \in \F,


\end{equation}


где $\al_i \in \R$, $i\in I$.





Минимальное $\sig$-замкнутое расширение назовем $\sig$-{\it


расширением}. Из смысла такого расширения видно, что


$$


\F = \bigcup_{k=0}^{+\infty} F_k,


$$


где $F_{k+1}=\Bigl\{\, \svij{i\in I}{} \al_i\, \max\limits_{j\in


J_i}\, f_j^i \mid f_j^i \in F_k,\; \al_i \in \R, \;


i \in I,\; j\in J_i, |I|<\pli,\; |J_i|< \pli \Bigr\}$.





На самом деле все функции из $\F$ могут быть преобразованы к


некоторому стандартному виду. В основе такого преобразования лежит


ряд тождеств, справедливых для произвольного набора функций:


\begin{gather}


\max_{j\in J}\, f_j + \max_{i\in I}\, g_i = \max_{(j,\,i) \in


J\times I} (f_j +g_i);\\


%


\sij{i\in I}{} \al_i\, \max_{j\in J_i}\, f_j^i =


\max_{s_i\in J_i}\, \sij{i\in I_+}{} \al_i\, f_{s_i}^i -


\max_{s_i\in J_i}\, \sij{i\in I_-}{} |\al_i|f_{s_i}^i,


\end{gather}


где $I_+ =\{ i\mid \al_i >0\}$, $I_- =\{ i\mid \al_i <0\}$;


\begin{gather}


\max_{j=1,\dotsc,m}\, f_j = \bigl[\, \max_{j=1,\dotsc,m-1}\, f_j


-f_m\bigr]^+ +f_m =


\bigl[f_m - \max_{j=1,\dotsc,m-1}\, f_j \bigr]^+ + \max_{j=1,\dotsc


,m-1}\, f_j; \\


%


\min_{i=1,\dotsc,n}\, f_i = - \bigl[\, \min_{i=1,\dotsc,n-1}\, f_i


-f_n\bigr]^+ + \min_{i=1,\dotsc,n-1}\, f_i =


-\bigl[f_n - \min_{i=1,\dotsc ,n-1}\, f_i \bigr]^+ +f_n;\\


%


\bigl[\, \max_{j\in J}\, f_j - \max_{i\in I}\, g_i \bigr]^+ =


\max_{(j,\,i) \in J\times I}\{f_j, g_i\} - \max_{i\in I}\, g_i;\\


%


\max_{j\in J}\, f_j - \max_{i\in I}\, g_i = \min_{i\in I}\,


\max_{j\in J}\, (f_j -g_i) = \max_{j\in J}\, \min_{i\in I}\, (f_j-g_i).


\end{gather}


Выписанные тождества носят общелогический характер и проверяются


непосредственно.





Наравне с $\F$ введем пространство


$$


\bold H = \bigcup_{i=0}^{+\infty}H_i,


$$


где $\bold H_0=\F_0$, $H_{k+1} =\Bigl\{\, \svij{i\in I}{}


\al_i\, f_i^+ \mid \al_i \in \R,\; f_i \in H_k, |I|< \pli \Bigr\}$.





Операция взятия положительной срезки ``$+$'' является частным


случаем $\sig$-операции, однако она потенциально (т.е. многократно


примененная) дает тот же класс функций $\F$. Имеет место





\begin{thms} % Теорема 1.1


$1)$ Все функции из $\F$ могут быть


приведены к любому из стандартных видов:


\begin{gather}


\max_{j\in J}\, f_j - \max_{i\in I}\, g_i,\\


%


\min_{i\in I}\, \max_{j\in J_i}\, f_j^i,\\


%


\max_{j\in J}\, \min_{i\in I_j}\, f_i^j,


\end{gather}


где $\{f_j, g_i, f_j^i \} \su


\F_0;$ {\em (1.9)} означает $F_k=F_1$ при $k>1$, следовательно,


$\F=F_1$.





$2)$ Классы $\F$ и $\bold H$ совпадают.





$3)$ Представления {\em (1.9)--(1.10)} эквивалентны.


\end{thms}





\begin{pf} 1) Соотношение (1.9) означает совпадение $F_k$ с


$F_1$ при $k>1$. На самом деле достаточно доказать $F_2=F_1$. В


силу (1.4) можно ограничиться преобразованием функции


$f(x)=\max\limits_{i\in I}\, f_i$ при $\{f_i\}_I \su F_1$ к


виду (1.9). Пусть $I=\{\ij{n}\}$. Доказательство можно вести


индукцией по $n$. Если $n=1$, то $f(x)=f_1(x)$ удовлетворяет


требуемому свойству представления (1.9). Пусть $n>1$. Так как


$f_i \in F_1$, то эти функции можно представить в виде


$$


f_i = \max_{j\in J_i}\, f_j^i - \max_{k\in I_i}\,g_k^i


$$


при $\{f_j^i, g_k^i \} \su \F_0$.


Имеем


$$


f(x) \overset{(1.5)}{=}{} \bigl[\, \max_{i=1,\dotsc,n-1}\, f_i -


f_n \bigr]^+ +f_n.


$$


По индукции


$\max\limits_{i=1,\dotsc,n-1}\, f_i =: \Bar{f} \in F_1$.


Так как $f= [\Bar{f} -f_n]^+ +f_n$, причем


$\{\Bar{f}, f_n \} \su F_1$, то по (1.4) $\Bar{f} -f_n \in F_1$, а


по (1.7) $[ \Bar{f} -f_n ]^+ \in F_1$, что (с


учетом $f_n \in F_1$) и дает $f\in F_1$. Итак, $F_2 =


F_1$, а вместе с тем и $\F=F_1$.





Представимость функций из $\F$ в виде (1.10) или


(1.11) следует из тождества (1.8).


\medskip





2) Докажем вначале включение $\bold H \su F_1$, т.е.


$H_k \su F_1$ $\fo\,k$. Если $f\in H_1$, то в


силу (1.4) $f\in F_1$. Следовательно, $H_1 \su F_1$.


Пусть $H_k \su F_1$, докажем $H_{k+1} \su F_1$. Произвольная


функция из $H_{k+1}$ имеет вид


$$


f = \sij{i\in I}{} \al_i\, f_i^+,\quad \{ f_i \} \su H_k\su F_1.


$$


Из этого следует, что $f$ представляется линейной комбинацией


функций дискретного максимума при образующих из $\F_0$, что с


учетом (1.4) и дает для $f$ требуемое представление (1.9).





Обратно, если $f\in \F$, т.е. $f$ имеет вид (1.9), то


показав, что функция дискретного максимума при образующих из $\F_0$


принадлежит $\bold H$, т.е. некоторому $H_k$, тем самым покажем и


$f\in \bold H$. Итак, пусть


$$


f = \max_{j\in J}\, f_j,\quad \{ f_j \} \su\F_0,\quad j=\ij{m}.


$$


Если $m=1$, то $f =f_1 \su \bold H_0$. Пусть $m>1$.


Если по индуктивному предположению $\Bar{f} =


\max\limits_{j=1,\dotsc,m-1}\, f_j \in H_k$, то в силу (1.5)


$f = [\Bar{f}-f_m]^+ +f_m \in H_{k+1}$, что и требовалось.


\medskip





3) Мы уже отметили, что представление функции в


форме (1.9) может быть переписано в форме (1.10) при


$f_j^i = f_j -g_i$ (см. (1.8)). Нужно убедиться в обратном.


Итак, пусть $f = \min\limits_{i\in I}\,


\max\limits_{j\in J_i}\, f_j^i$. Если $I=\{\ij{n}\}$ и


$n=1$, то $f=\max\limits_{j\in J_1}\, f_j^1 \in F_1$. При


$n>1$ доказательство, как это и делалось выше, можно провести


индукцией по $n$. По (1.6)


$$


f = \bigl[\, \min_{i=1,\dotsc,n-1}\, \max_{j\in J_i}\, f_j^i


- \max_{j\in J_n}\, f_j^n \bigr]^+ + \max_{j\in J_n}\,f_j^n.


$$


По индуктивному предположению


$$


\Bar{f} = \min_{i=1,\dotsc,n-1}\, \max_{j\in J_i}\,


f_j^i \in F_1,


$$


т.е. $\Bar{f}$ может быть представлено в форме (1.9), но


тогда с использованием преобразований (1.3), (1.4) и


(1.7) функция $f$ приводится к виду (1.9), что и


требовалось.





Эквивалентность представлений (1.11) и (1.9)


доказывается аналогично.


\end{pf}





Функции $f$ из $\F$ будем называть $\sig$-{\it функциями} или


$\sig$-{\it кусочными функциями}. Если $\F_0$ --- \linebreak


пространство


линейных (аффинных) функций, то $\F$ --- пространство


кусочно-линейных функций.





\section{Кусочно-линейные функции}





Кусочно-линейные функции --- это $\sig$-функции в


ситуации, когда $\F_0$ --- пространство линейных


(аффинных) функций. К определению кусочно-линейных (в


дальнейшем --- $k$-{\it функций}) можно подойти двояко: либо так,


как это только что сделано, либо исходя из некоторой аксиоматики,


идентифицирующей такие функции. Остановимся на втором подходе.





Пусть имеются конечные совокупности многогранников $\{ M_j \}_J$


и собственных линейных функций $\{l_j(x) \}_J$. Будем


говорить, что система $\{ M_j, l_j(x) \}$ задает


однозначную {\it кусочно-линейную функцию} $l(x)$ на $\X$, если:


\begin{itemize}


\item[1)] $\bigcup\limits_{j \in J}\, M_j =\X,\quad M_i^0


\bigcap M_j^0 = \emp\; \text{ при }\; i \ne j$;


\item[2)] $l(x) \equiv l_j (x)\quad \fo\, x \in


M_j,\quad \fo\, j \in J$.


\end{itemize}


Здесь $M_j^0$ --- {\it алгебраическая} внутренность 


многогранника $M_j$, т.е. $y \in M_j^0 \,\Llra\, y +t s \in


M_j$ при любом $s \in\X$ и достаточно малом $t > 0$. Термином


{\it многогранник}, как и ранее, обозначается множество, задаваемое


конечной системой собственных линейных неравенств


$$


(f_j, x) -\al_j \le 0,\quad j\in J.


$$


В данном определении некоторые из многогранников $M_j$ или~$M_j^0$ могут


быть и пустыми.





Будем использовать обозначения: $\bold L_0$ --- пространство


аффинных функций, $\bold L$ --- пространство $k$-функций, определенных


внешним образом в силу свойств 1) и 2). Действительно, класс


функций $\bold L$ совпадает с классом $\F$ из предыдущего параграфа,


который строится из $\F_0 =\bold L_0$, т.е. $\bold L$ --- это


минимальное алгебраическое расширение пространства аффинных функций,


замкнутое относительно операций дискретного максимума \cite{Plo}.


Следовательно, представление (1.9) (а также каждое


из (1.10) и (1.11)) является универсальным


представлением кусочно-линейных функций ($k$-функций).





Для унификации и упрощения записи $k$-функций, систем неравенств


из $k$-функций, задач кусочно-линейного программирования и т.д.


будем в дальнейшем полагать $\X=\R^n$, тогда


\begin{gather*}


\bold L_0 = \{ l(x) = (a, x) - \al \mid a \in \R^n,\; \al \in \R \},\\


\bold L = \bigl\{ \max_{j \in J}\, l_j(x) - \max_{i \in I}\,


h_i (x) \mid \{ l_j, h_i \} \su \bold L_0,\;


|J|< +\infty,\; |I|< +\infty\bigr\}.


\end{gather*}





Введем обозначения:\


$$


|z|_{\max} = \max_{(i)}\, z_i,\quad |z|_{\min} =\min_{(i)}\, z_i,


$$


где $z$ --- вектор конечномерного пространства.


Если $Ax -b =[\nk{l_1(x)}{l_m(x)}]^T$ --- вектор


аффинных функций, то в соответствии с введенным обозначением будем


иметь


$$


| Ax - b\, |_{\max} = \max_{(i)}\, l_i(x).


$$


Представления кусочно-линейных функций в виде (1.9)--(1.11)


принимают унифицированную форму


\begin{gather}


|Ax -b|_{\max} - |Bx -d|_{\max},\\


%


\min_i |A_i x -b^i|_{\max},\\


%


\max_{(j)}\, | A_j x -b^j|_{\min}.


\end{gather}





Отметим также свойства функций дискретного максимума


\begin{gather*}


|z|_{\max} = - |z|_{\min};\\


\text{при } \al > 0:\; |\al\, z|_{\max} = \al\,


|z|_{\max};\qquad \text{при } \al < 0:\; |\al\,


z|_{\max} = \al\, |z|_{\min}.


\end{gather*}





\section{Системы кусочно-линейных неравенств и их


геометрическая интерпретация}





Конечная система $k$-функций может быть записана в виде


\begin{equation}


\min_{(j)}\, |A_j^t x -b_t^j|_{\max} \le 0,\quad t=\ij{T},


\end{equation}


или с помощью одного неравенства


$$


\sum_{t=1}^T \min_{(j)}\, |A_j^t x -b_t^j|_{\max} \le 0,


$$


или (на основе теоремы 1.1) в виде


\begin{equation}


\min_{j=1,\dotsc, m}\, |A_jx -b^j|_{\max} \le 0.


\end{equation}


Это задание будем считать одним из стандартных. Другим


стандартным видом является


\begin{equation}


|Ax - b|_{\max} - |Bx - d|_{\max} \le 0


\end{equation}


(см. (2.1)). Итак, произвольная конечная система


кусочно-линейных неравенств может быть приведена к \


любому из стандартных видов (3.1)--(3.3).





Рассмотрим представление системы в виде (3.2).


Положим $M_j =\{ x \mid A_j x \le b^j \}$. Тогда


множеством решений неравенства (3.2) будет $M =


\bigcup\limits_{j=1}^m\, M_j$. С другой стороны, если $M$ ---


произвольное полиэдральное множество, пусть из $\R^n$, т.е. $M =


\bigcup\limits_{j=1}^m\, M_j$ и $M_j$ --- многогранники,


которые задаются конечными системами линейных неравенств ($M_j


= \{ x \mid A_j x \le b^j \}$), то $M$ является множеством решений


неравенства (3.2).





С этой же точки зрения посмотрим на неравенство (3.3).


Пусть


\begin{gather*}


Ax -b = [\nk{l_1 (x)}{l_m (x)} \,]^T,\qquad


Bx -d = [\nk{s_1 (x)}{s_k (x)} \,]^T,\\


\{ l_j(x), s_i(x) \}_{1,\,1}^{m,\,k} \su \bold L_0


\end{gather*}


(\,$\bold L_0$ --- пространство аффинных функций). Положим


$$


M_i = \{ x \mid l_j(x) \le s_i(x),\; j=\ij{m} \}.


$$


Тогда, как легко видеть, $\bigcup\limits_{i=1}^s\, M_i$


совпадает с множеством решений неравенства (3.3). Тем самым


для неравенства (3.3) указаны те многогранники $M_i$,


объединение которых и дает множество решений неравенства (3.3).


С другой стороны, если полиэдральное множество $M$ задано тем или


иным образом, например, как в предыдущем случае, --- совокупностью


систем линейных неравенств, каждая из которых задает выпуклую


многогранную компоненту множества $M$, то требуется цепочка тех или


иных преобразований типа (1.3)--(1.8), приводящая к


заданию множества $M$ одним неравенством вида (3.3).





Наконец, рассмотрим систему кусочно-линейных неравенств в


форме (3.1), формально более сложную, чем (3.2) или


(3.3). Положим


$$


M_j^t:= \{ x \mid A_j^t x \le b_t^j \},\qquad


M_t:= \bigcup_{(j)}\, M_j^t,\qquad M:= \bigcap_{(t)}\, M_t.


$$


Множество $M_t$ --- это множество решений $t$-го неравенства в


системе (3.1), так что $M$ --- множество решений всей системы.





Материал, изложенный в \S\,1--\S\,3, устанавливает


способы конструктивного соответствия между полиэдральными множествами


и их аналитическим заданием. Хотя сопутствующая этому алгебра


преобразований может быть достаточно громоздкой, тем не менее логика


этих преобразований проста и может быть в реальных прикладных


ситуациях поручена компьютеру.





\section{Задача кусочно-линейного


программирования}





4.1. {\it Предварительные замечания}.


Произвольной задаче кусочно-линейного программирования, т.е.


задаче поиска экстремума $k$-функции при ограничениях в форме


конечной системы неравенств с $k$-функциями в левых частях, можно


придать универсальный и простой вид


\begin{equation}


P: \max \{(c, x)\bigl|{}\bigr. \min_{j=1,\dotsc,m}\,


|A_j x -b^j|_{\max} \le 0,\; x \ge 0 \}.


\end{equation}


Действительно, о способе сведения системы $k$-неравенств к


одному $k$-неравенству уже говорилось. Если же $f(x)$ ---


произвольная оптимизируемая, пусть максимизируемая, $k$-функция при


одном $k$-неравенстве, пусть $g (x) \le 0$ (возможно с $x \ge


0$), то переписав задачу $\max\, \{ f(x)|g (x) \le 0 \}$ в


виде


$$


\max\, \{ t \mid g(x) \le 0,\; f(x) \ge t \}


$$


и преобразовав систему из двух $k$-неравенств к одному


$k$-неравенству, получим задачу поиска максимума линейной функции при


одном ограничении в форме $k$-неравенства. В (4.1) ограничение


$x \ge 0$ выделено для целей обеспечения симметрии в ряде


аналитических конструкций, рассматриваемых ниже.





Итак, объектом рассмотрения примем задачу (4.1). Введем


{\it частичную} задачу


\begin{equation}


L_j: \max\, \{ (c, x) \mid A_j x \le b^j,\; x \ge 0 \}.


\end{equation}


Связь между задачей (4.1) и $L_j$ проста:


\begin{equation}


\opt (4.1) = \max_{(j: M_j \ne \emp)}\,\opt L_j,


\end{equation}


где $M_j = \{ x \ge 0 \mid A_j x \le b^j \}$. В (4.2)


для некоторых $j$ множества $M_j = \{ x \ge 0 \mid A_j x \le b^j \}$


могут быть пустыми при свойстве разрешимости исходной


задачи (4.1), т.е. произвольная \linebreak


$k$-задача, будучи


преобразована к виду (4.1), распадается на конечное число задач


линейного программирования, решив которые, тем


самым находим решение и исходной задачи. Поскольку


конструктивизм соответствующих


преобразований обеспечен, то формально решение произвольной


задачи кусочно-линейного программирования сводится к использованию


отмеченного конструктивизма и некоторого метода (напр.,


симплекс-метода) решения задачи ЛП.


\medskip





4.2. {\it Условия разрешимости $k$-задачи}.


Для $k$-задач выполняется ряд свойств и теорем, формулируемых в


рамках теории линейного программирования. Некоторые из них являются


следствиями теорем из ЛП, некоторые же требуют своих доказательств,


по крайней мере, уточнений. Не требует самостоятельного доказательства





\begin{thms}[{\series{bx}\normalshape\selectfont{}о достижимости}]


Если


$$


\sup\, \{ (c, x) \bigl|{}\bigr. \min_{(j)}\,


|A_j x - b^j|_{\max} \le 0,\; x \ge 0 \} < \pli,


$$


то операция $\sup$ в этой задаче достижима.


\end{thms}





Выпишем задачу, двойственную к~$L_j$:


\begin{equation}


L_j^*: \min\, \{ (b^j, u^j) \mid A_j^T u^j \ge c,\; u^j \ge 0\}.


\end{equation}


Пусть $M_j^* = \{ u^j \ge 0 \mid A_j^T u^j \ge c \}$.





\begin{thms}


Задача {\em (4.1)} разрешима тогда и только тогда, когда


$$


M\, =\, \bigcup_{j=1}^m\, M_j \ne \emp\qquad \text{и} \qquad


M_j \ne \emp \, \Lra\, M_j^* \ne \emp.


$$


\end{thms}





\begin{pf} Поскольку условия $M_j \ne \emp$ и $M_j^* \ne \emp$


являются необходимыми и достаточными для разрешимости задачи $L_j$,


то $\max\limits_{j: M_j \ne \emp}\, \opt L_j$


конечен и является значением $\opt P$. Обратно, если задача $P$


разрешима, то все задачи $L_j$ c $M_j \ne \emp$ разрешимы


(т.е. ситуации $\opt L_j = \pli$ исключаются), а тогда в


соответствии с теоремой двойственности в ЛП $M_j^* \ne \emp$, что


и требовалось.


\end{pf}





\begin{notes}


Другой вариант формулировки теоремы 4.2 состоит в


эквивалентности


$$


(P \text{ разрешима }) \,\Llra\, (M \ne \emp \,\&\, M_j \ne


\emp \,\Lra\, L_j \text{ разрешима}).


$$


\end{notes}





4.3. {\it Двойственность}.


Исходную $k$-задачу возьмем в форме (4.1), дополнив


ее предположением (впрочем, несущественным): размерность векторов


$A_j x -b^j$ одинакова, т.е. число неравенств в каждой из систем


$A_j x \le b^j$ одно и то же. Введенное условие позволяет


двойственную переменную для $L_j$, т.е. переменную $u^j$ в


задаче (4.4) обозначать одним символом $u$. В соответствии со


сказанным задачу (4.4) перепишем в виде


\begin{equation}


\min\, \{ (b^j, u) \mid A_j^T u \ge c,\; u \ge 0 \},\quad


j=\ij{m}.


\end{equation}


Задачу


\begin{equation}


P^*: \max_{j: M_j^* \ne \emp}\, \min\, \{ (b^j, u) \mid


A_j^T u \ge c,\; u \ge 0 \}


\end{equation}


будем рассматривать как двойственную к $P$. По


отношению к форме задания $P$ задача $P^*$ не


симметрична, но если $P$ переписать в эквивалентном виде


\begin{equation}


\max_{j\in \overline{1,m}}\, \min\, \{ (c, x) \mid A_j x \ge


b^j,\; x \ge 0 \},


\end{equation}


то (4.6) и (4.7) принимают симметричный вид.





В отношении задачи (4.6) справедлива





\begin{thms}


Задача {\em (4.6)} разрешима тогда и только тогда, когда


\begin{equation}


\exists\, j \in \overline{1,\, m}: M_j^* \ne \emp \,\& \, M_j \ne \emp.


\end{equation}


\end{thms}





\begin{pf} Действительно, если $J:= \{ j \mid M_j \ne \emp, M_j^*


\ne \emp \}$, то для $j \notin J: \inf\, \{ (b^j, u) \mid x


\in M_j^* \} = -\infty$, поэтому в (4.6) максимум можно


брать только по $j \in J$. Этим и обеспечивается разрешимость


задачи $P^*$. Необходимость условий (4.8) очевидна: если $P^*$


разрешима, то хотя бы для одного $j'$ задача $L_j^*$ разрешима, что


эквивалентно (4.8).


\end{pf}





Из теорем 4.2 и 4.3, а также теоремы


двойственности в ЛП вытекает





\begin{thms}


Если задача {\em (4.1)} разрешима, то {\em (4.6)} также


разрешима, при этом их оптимальные значения совпадают.


\end{thms}





В ситуации задач $P$ и $P^*$ свойство одновременной их


разрешимости или неразрешимости в отличие от ЛП теряется.


Принимая во внимание возможность несобственных оптимальных значений,


запишем эти задачи в виде


$$


P: \sup\limits_{(j)}\, \inf\limits_{x \in M_j}\, (c,\


x),\qquad


P^*: \sup\limits_{(j)}\, \inf\limits_{u \in M_j^*}\, (b^j, u).


$$


Условия одновременной разрешимости $P$ и $P^*$ определяются


теоремами 4.2 и 4.3, но возможна и ситуация:


$P$ не разрешима, $P^*$ разрешима. Это соответствует тому,


что при разрешимости~$P^*\;$ $\exists\, j_0:M_{j_0} \ne \emp\;


\&\; M_{j_0}^* = \emp$, т.е. задача $L_{j_0}$ --- несобственная 2-го


рода. Так как в рассматриваемой ситуации


$\sup\limits_{x \in M_{j_0}}\, (c, x) = + \infty$, то $\opt P = + \infty$,


т.е. $P$ не разрешима. Одновременная неразрешимость


задач $P$ и $P^*$ реализуется на паре задач линейного


программирования $L$ и $L^*$, являющихся несобственными 3-го


рода.





\section{Задача о седловой точке


дизъюнктивной функции Лагранжа}





Пусть $\{ M_j \}_1^m \su \R^n$ и $f(x)$ ---


произвольная функция, заданная на $\R^n$. Запишем задачи


\begin{equation}


\begin{split}


P_{\cap}&: \max\, \bigl\{ f(x) \mid x \in \bigcap_{(j)}\, M_j \bigr\},\\


%


P_{\cup}&: \max\, \bigl\{ f(x) \mid x \in \bigcup_{(j)}\, M_j \bigr\}.


\end{split}


\end{equation}


Первую из них назовем задачей в {\it конъюнктивной}


постановке, вторую --- в {\it дизъюнктивной}. Первая из постановок


обычна для традиционных постановок задач


математического программирования. Вторая из них


отражает ситуацию кусочно-линейной задачи (4.1), которую можно


переписать в виде


$$


\max\, \bigl\{ (c, x) \mid x \in \bigcup_{(j)} M_j \bigr\}


$$


при $M_j =\{ x \mid A_j (x) \le b^j,\; x \ge 0 \}$, $j=\ij{m}$.





Рассмотрим задачу (5.1) при условиях


$$


M_j =\{ x \mid F_j (x) \le 0,\; x \ge 0 \},\quad


F_j: \R^n \,\to\, \R^{m_j},\quad j=\ij{m}.


$$


Зафиксируем схему соответствия задачам


$P_{\cap}$ и $P_{\cup}$ их функций Лагранжа


\begin{equation}


\begin{split}


P_{\cap} & \to f(x) -\sum_{j = 1}^m (u_j, F_j(x))=\varPhi_{\cap}(x,u),\\


P_{\cup} & \to f(x) -\min_{(j)}\, (u_j, F_j (x))=\varPhi_{\cup}(x,u)


\end{split}


\end{equation}


--- {\it дизъюнктивная функция} Лагранжа для $P_{\,\cup}$. 





Задачу (5.1) назовем {\it вполне регулярной}, если


каждая из задач


\begin{equation}


\max\, \{ f(x) \mid F_j (x) \le 0,\; x \ge 0 \}\tag{$5.1_j$}


\end{equation}


разрешима. Эквивалентом такой регулярности является


разрешимость задачи и $M_j \ne \emp$, $j=\ij{m}$.





Используем в дальнейшем следующее обозначение: если экстремальной


задаче приписан номер $t$, то ее решение или множество решений


будет обозначаться через $\arg (t)$, $\Arg(t)$ соответственно,


а оптимальное значение --- через $\opt(t)$.





Пусть $\varPhi_j (x, u_j) = f(x) - (u_j, F_j (x))$ ---


функция Лагранжа для (5.1)$_j$, $x \in \R^n$, $u_j


\in \R^{m_j}$. Справедлива хорошо известная (напр., \cite{AHU})





\begin{lems}


Если точка $[ \Bar{x}, \Bar{u} ] \ge 0$ является седловой


для функции Лагранжа $\varPhi (x, u) = f(x) - (u, F(x))$,


поставленной в соответствие задаче


\begin{equation}


\max\, \{ f(x) \mid F(x) \le 0,\; x \ge 0 \},


\end{equation}


то $(\Bar{u}, F(\Bar{x})) =0$ и $\bar{x} \in \Arg (5.3)$.


\end{lems}





\begin{lems}


Пусть каждая из функций $\varPhi_j (x, u_j)$ обладает


седлом $[ \Bar{x}_j, \Bar{u}_j ] \ge 0$, $j=\ij{m}$.


Если $\Bar{x}$ --- такое значение $\Bar{x}_j$, на котором достигается


$\max\limits_{(j)}\, f(\Bar{x}_j)$, то


\begin{equation}


\varPhi (x, \Bar{u})\le f(\Bar{x}) \ \; \forall x \ge 0,


\end{equation}


где $\varPhi (x, \Bar{u}) = f(x) - \min\limits_{(j)}\,


(\Bar{u}_j, F_j (x))$.


\end{lems}





{\bf Доказательство.} Согласно лемме 5.1\; $(\Bar{u}_j, F_j


(\Bar{x}_j)) = 0$, $j=\ij{m}$, и $f(x) - (\Bar{u}_j, F_j (x))


\le f(\Bar{x}_j)$ $(\le f(\Bar{x}))$ $\fo\, x \ge 0$.


Отсюда


$$


\max_{(j)}\, [f(x) -(\Bar{u}_j, F_j (x))] \le f(x)


-\min_{(j)}\, (\Bar{u}_j, F_j(x)) = \varPhi (x, \Bar{u}) \le


f(\Bar{x})\; \fo\, x \ge 0.\quad\square


$$





\begin{lems}


Если $[ \Bar{x}, \Bar{u} ] \ge 0$ --- седловая точка


для $\varPhi (x, u)$, то $\min\limits_{(j)}\, (\Bar{u}_j,


F_j(\Bar{x})) =0$ и $\Bar{x} \in \Arg(5.1)$.


\end{lems}





\begin{pf} В соответствии с определением седловой точки


$$


\varPhi (x, \Bar{u})


\underset{\forall x\ge 0}{\le}{}


\varPhi (\Bar{x}, \Bar{u})


\underset{\forall u\ge 0}{\le}{}


\varPhi (\Bar{x}, u),


$$


или


\begin{gather}


f(x) -\min_{(j)}\, (\Bar{u}_j, F_j(x))


\underset{\forall x\ge 0}{\le}{}


f(\Bar{x}) -\min_{(j)} (\Bar{u}_j, F_j(\Bar{x})),\\


%


f(\bar{x}) -\min_{(j)}\, (\bar{u}_j, F_j (\bar{x}))


\underset{\forall u\ge 0}{\le}{}


f(\Bar{x}) -\min_{(j)}\, (u_j, F_j (\Bar{x})).


\end{gather}


Соотношение (5.6) перепишем в виде


\begin{equation}


\min_{(j)}\, (\Bar{u}_j, F_j (\Bar{x}))


\underset{\forall u\ge 0}{\le}{}


\min_{(j)}\, (u_j, F_j (\Bar{x})).


\end{equation}


Докажем, во-первых, $\Bar{x} \in M = \bigcup\limits_{(j)}\,


M_j$, т.е. $j_0: F_{j_0} (\Bar{x}) \le 0$. Если бы


$F_j (\Bar{x}) \not\le 0$ $\fo\, j$, то за счет выбора $u \ge 0$


можно было бы значение правой части в (5.7) сделать как угодно


большим, что будет противоречить соотношению (5.7). Доказанное


дает $\al:= \min\limits_{(j)}\, (\Bar{u}_j, F_j (\Bar{x})) \le


0$. На самом деле, $\al = 0$. Если бы $\al < 0$, то


соотношение (5.7) при $u = 0$ порождало бы противоречивое


неравенство $0 > - |\al|\ge 0$. Итак, соотношение


$\min\limits_{(j)}\, (\Bar{u}_j, F_j (\Bar{x}))$ доказано.





Необходимо доказать оптимальность вектора~$\Bar{x}$ 


для (5.1), т.е. $f(x) \le f(\Bar{x})$ $\fo\, x \in


M$. Обратимся к соотношению (5.5), которое можно теперь


переписать в виде


\begin{equation}


f(\Bar{x}) - f(x) \ge -\min_{(j)}\, (\Bar{u}_j, F_j (x)).


\end{equation}


Если $x \in M$, т.е. $x \in M_{j_0}$ при некотором $j_0$,


то $\min\limits_{(j)}\, (\Bar{u}_j, F_j (x)) \le 0$, что в


силу (5.8) дает $f(\Bar{x}) -f(x) \ge 0$ $\fo\, x \in


\bigcup\limits_{(j)}\, M_j$.


\end{pf}





\begin{lems}


Пусть задача {\em (5.1)}, т.е.


$$


L_{\cup}: \max\, \bigl\{ (c, x) \mid x \in \bigcup_{j=1}^m\, M_j \bigr\},


$$


при $M_j = \{ x \mid A_j x \le b^j,\; x \ge 0 \}$


разрешима и $M_j \ne \emp$ $\fo\,j$ {\em (}т.е. вполне регулярна{\em )}.


Тогда функция $L_{\cup}(x, u) =(c, x) -\min\limits_{(j)}\, (u_j,\


A_j x - b^j)$ обладает седлом $[ \Bar{x}, \Bar{u} ]$, причем в


качестве $\Bar{x}$ и $\Bar{u} = [ \nk{\Bar{u}_1}{\Bar{u}_m} ]$


выступают


$$


\Bar{x} =\arg\, \max_{(j)}\, (c, \Bar{x}_j),\quad


\Bar{x}_j \in \Arg\, \max_{x \in M_j}\, (c, x),\quad


\Bar{u}_j \in \Arg\, \min_{u \in M^*}\, (b^j, u),\quad j=\ij{m}.


$$


\end{lems}





\begin{pf} Если показать, что $\al:= \min\limits_{(j)}\,


(\Bar{u}_j, A_j \Bar{x} -b^j) = 0$, то в силу леммы 5.2


левое неравенство в определении седла для $L_{\cup} (x, u)$ будет


выполнено. Так как вектор $\Bar{x} \ge 0$ удовлетворяет одной из


систем $A_j x \le b^j$, то $\al \le 0$. Покажем $\al \ge 0$. Имеем


$$


\min_{(j)}\, (\Bar{u}_j, A_j \Bar{x} -b^j) =


\min_{(j)}\, [ -(b^j, \Bar{u}_j) +(c, \Bar{x}) + (A_j^T


\Bar{u}_j - c, \Bar{x}) ] \ge \min_j\, [ (c, \Bar{x}) - (c,


\Bar{x}_j) ] \ge 0.


$$


Выше мы воспользовались соотношениями: $A_j^T \Bar{u}_j - c


\ge 0$; $(b^j, \Bar{u}_j) = (c, \Bar{x}_j)$ --- по теореме


двойственности в ЛП; $(c, \Bar{x}) \ge (c, \Bar{x}_j)$ $\fo\, j$.





Остается показать выполнимость правого неравенства в определении


седловой точки для $L_{\cup} (x, u)$, т.е.


$$


L_{\cup} (\Bar{x}, \Bar{u})


\underset{\forall u\ge 0}{\le}{}


L_{\cup} (\Bar{x}, u).


$$


С учетом $\al = 0$ выписанное неравенство принимает вид


\begin{equation}


0 \underset{\forall u_j\ge 0}{\le}{}


-\min_{(j)}\, (u_j, A_j \Bar{x} - b^j).


\end{equation}


Так как $\fo\, j_0$ $A_{j_0} \Bar{x} -b^j \le 0$, то при


любом $u \le 0\;$ $\min\limits_{(j)}\, (u_j, A_j \Bar{x} -b^j) \le


0$, следовательно, (5.9) верно.


\end{pf}





Из лемм 5.3 и 5.4 вытекает





\begin{thms}


Если системы $A_j x \le b^j$, $x \ge 0$, $j=\ij{m}$,


совместны, то задача {\em (5.1)} разрешима тогда и только тогда,


когда ее дизъюнктивная функция Лагранжа


$$


L_{\cup} (x, u) = (c, x) -\min_{(j)}\, (u_j, A_j x - b^j)


$$


обладает седлом $[ \Bar{x}, \Bar{u} ] \ge 0$, при этом


\begin{itemize}


\item[1)] если $L_{\cup}$ разрешима и $\Bar{x}_j \in


\Arg L_j$, $\Bar{u}_j \in \Arg L_j^*$, $\Bar{x} =


\arg\, \max\limits_{(j)}\, (c, \Bar{x}_j)$, то $[ \Bar{x},


\Bar{u} ]$ --- седло$;$


\item[2)] если $[ \Bar{x}, \Bar{u} ]$ --- седло для $L_{\cup}


(x, u)$, то $\{ \Bar{x}, \Bar{u}_j \}$ удовлетворяют всем


соотношениям из {\em 1)}.


\end{itemize}


\end{thms}





Существование седловой точки $[ \Bar{x}, \Bar{u} ] \ge 0$


функции $\varPhi (x, u)$, пусть в форме (5.5)--(5.6),


эквивалентно выполнимости соотношения


\begin{equation}


\max_{x \ge 0}\, \min_{u \ge 0}\, \varPhi (x, u) = \min_{u


\ge 0}\, \max_{x \ge 0}\, \varPhi (x, u) = \varPhi (\Bar{x},


\Bar{u}).


\end{equation}





В игровой интерпретации двойственности в


математическом программировании 


утверждения типа теоремы 5.1 естественно давать через


соотношение (5.1), а именно, справедлива





\begin{thms}


Пусть задача {\em (5.1)} разрешима и $M_j \ne \emp$ $\fo\,


j$. Тогда для $\Bar{x} \in \Arg (5.1)$ существуют


$\Bar{u}_j \ge 0$, $j=\ij{m}$, такие, что вектор $[ \Bar{x},


\Bar{u} ]$ будет удовлетворять соотношению {\em (5.10)} при


$$


\varPhi (x, u) = L_{\cup} (x, u) = (c, u) - \min_{(j)}\,


(\Bar{u}_j, A_j x - b^j).


$$


\end{thms}





\section{Метод точных штрафных функций для задачи\protect\\


кусочно-линейного программирования}





Общую задачу кусочно-линейного программирования в


канонической постановке (4.1), т.е.


\begin{equation}


P_{\cup}: \max\, \{ (c, x) \bigl|{}\bigr. \min_{j=1,\dotsc,m}\,


|A_j x -b^j|_{\,\max} \le 0,\; x \ge 0 \},


\end{equation}


можно привести к эквивалентной задаче этого же типа,


но со снятием основного ограничения в (6.1). Поставим задаче


$L_{\cup}$ в соответствие $k$-задачу


\begin{equation}


\sup_{x \ge 0}\, \bigl[(c, x) -\min_{(j)}\, (R_j, (A_j x -b^j)^+\, )\bigr],


\end{equation}


где $R_j$ --- неотрицательный векторный параметр


размерности~$m_j$, т.е $m_j$ --- число неравенств в системе $A_j


x - b^j \le 0$.





Как и в предыдущем параграфе, будем использовать следующие


обозначения: $L_j$ --- это задача


$\max\, \{ (c, x) \mid A_j x \le b^j,\; x \ge 0 \}$,


$L_j^*$ --- двойственная к ней,


$\Bar{u}_j =\arg L_j^*$, $\Bar{u} = [


\nk{\Bar{u}_1}{\Bar{u}_m} ]$, $M_j = \{ x \ge 0 \mid A_j x \le


b^j \}$.





\begin{thms}


Пусть задача {\em (6.1)} разрешима, $M_j \ne


\emp$ $\fo\, j;$ $\Bar{u}_j \in


\Arg L_j^{\ast}$. Если $R_j \ge R_0\, \Bar{u}_j$, $R_0 >


1$, то оптимальные значения и оптимальные множества задач


{\em (6.1)} и {\em (6.2)} совпадают, т.е.


\begin{align}


\opt (6.1) & =\opt (6.2),\\


%


\Arg (6.1) &=\Arg (6.2).


\end{align}


\end{thms}





\begin{pf} Обозначим целевую функцию в (6.2)


через $\varPhi_R (x)$, а вычитаемую из $(c, x)$ часть --- через


$\varPhi_0 (x)$. Докажем вначале равенство (6.3). Для


$\Bar{x} \in\Arg(6.1)$ получаем $\varPhi_R (\Bar{x}) =


(c, \Bar{x})\, =\opt(6.1)$, следовательно,


$\opt(6.2)=\sup\limits_{x \ge 0}\, \varPhi_R (x) \ge\opt(6.1)$.





Докажем обратное неравенство. По леммам 5.4  и 5.3


$$


(c, x) -\min_{(j)}\, (\Bar{u}_j, A_j x -b^j) \le (c,\


\Bar{x})\quad \fo\, x \ge 0.


$$


С учетом этого неравенства оценим $\varPhi_R(x)$ для $x \ge 0$


\begin{multline}


\varPhi_R(x) \le (c, \Bar{x}) +\min\limits_{(j)}\,


(\Bar{u}_j, \;A_j x - b^j) -\varPhi_0 (x) \le \\


%


\le \opt(6.1) + \min\limits_{(j)}\, (\Bar{u}_j,\;


(A_j x - b^j)^+ ) - \varPhi_0(x) \le 


\opt(6.1)+ \frac{1}{R_0} (R_j,\;


(A_j x - b^j)^+ ) - \varPhi_0 (x) = \\


%


= \opt(6.1) - \frac{R_0 - 1}{R_0}\,


\min\limits_{(j)}\, (R_j,\; (A_j x - b^j)^+ ) \le


\opt(6.1).


\end{multline}


Отсюда $\sup\limits_{x \ge 0}\, \varPhi_R (x)


\le \opt(6.1)$. Таким образом, равенство (6.3)


доказано. Из него, в частности, следует включение


$\Arg(6.1) \su \Arg(6.2)$, что позволяет, на самом


деле, в задаче (6.2) вместо $\sup$ писать $\max$.


Докажем обратное включение. Пусть $\Bar{x} \in\Arg(6.2)$.


Согласно (6.5) имеем


$$


\opt(6.1) = \varPhi_R (\Bar{x}) \le


\opt(6.1) - \frac{R_0 - 1}{R_0}\,


\min_{(j)}\, (R_j,\; (A_j \Bar{x} - b^j)^+ ).


$$


Отсюда вытекает


$\min_{(j)}\,(R_j,\; (A_j \Bar{x} - b^j)^+ \,) = 0,$


а потому $\exists\, j_0: (R_{j_0},\; (A_{j_0} \Bar{x} -


b^{j_0})^+ ) = 0$, что ввиду $R_{j_0} > 0$ дает $A_{j_0}


\Bar{x} \le b^{j_0}$. Следовательно, $\Bar{x} \in M_{j_0} \su M


= \bigcup\limits_{j=1}^m\, M_j$. Допустимость вектора $\Bar{x}$ для


задачи (6.1) вместе с $(c, \Bar{x}) =\opt(6.1)$


означает $\Bar{x} \in\Arg(6.1)$, а потому и


$\Arg(6.2) \su \Arg(6.1)$. Доказательство


равенства (6.4) также завершено.


\end{pf}





Конструкция доказательства может быть повторена для более общей


ситуации, а именно для задачи (5.2) и эквивалентной редукции ее


к задаче


\begin{equation}


\sup_{x \ge 0}\, \bigl[ f(x) - \min_{(j)}\, (R_j, F_j^+ (x)) \bigr].


\end{equation}


Справедлива





\begin{thms}


Пусть каждая из задач $\max\, \{ f(x) \mid F_j (x) \le 0,\; x


\ge 0 \}$ обладает седлом $[ \Bar{x}_j, \Bar{u}_j ]$. Тогда при


$R_j > R_0\,\Bar{u}_j$, $R_0 > 1$, задачи {\em (5.2)} и


{\em (6.6)} эквивалентны в смысле совпадения их оптимальных


значений и оптимальных множеств.


\end{thms}





Действительно, в соответствии с леммой 5.3 можно


выписать неравенство


$$


f(x) \le f(\Bar{x}) + \min_{(j)}\, (\Bar{u}_j, F_j(x))\quad


\fo\, x \ge 0,


$$


а далее все повторить в соответствии с выкладками (6.5).





\section{Вопросы полиэдральной разделимости}





Задача разделимости непересекающихся множеств $M$ и


$N$ из некоторого пространства $\bold X$ с помощью функции $f(x)$ из


фиксированного класса ${\cal F}_0$ называется {\it дискриминацией


множеств}. Функция $f(x)$ в этом случае называется


{\it дискриминантной}. Разделимость состоит в выполнимости


неравенств


\begin{equation}


f(x) > 0\quad \fo\, x \in M;\qquad f(y) < 0\quad \fo\, y \in N.


\end{equation}


Можно говорить и о нестрогой разделимости, когда в (7.1)


допускаются нестрогие неравенства.





Если $M$ и $N$ --- выпуклые многогранники, а ${\cal F}_0$


--- класс аффинных функций, то говорят о задаче {\it линейной


дискриминации}. Задача линейной дискриминации является одной из


важнейших задач в распознавании образов (РО). В качестве одной из


базовых формальных моделей РО является система строгих линейных


однородных неравенств. Убедимся в этом.





Пусть ${\cal G}$ и ${\cal L}$ --- некоторые {\it образы}, $A


= \{a_j\} \su\R^n$ и $B = \{b_i\} \su\R^n$ ---


формализованные выборки представителей этих образов. Если


$\{f_s(x)\}_1^k$ --- некоторый базовый набор функций (вообще говоря,


произвольный), то разделяющую функцию можно искать в виде $f (x) =


\svij{s=1}{k}z_s\,f_s(x)$, где $\{z_s\}$ --- числовые


коэффициенты. Свойство разделимости состоит в выполнимости системы


неравенств


$$


\sij{s = 1}{k}z_s\, f_s (a_j) > 0\; \fo\, j\,;\qquad


\sij{s=1}{k} z_s\, f_s (b_i) < 0\; \fo\, i,


$$


или в матричном виде


\begin{equation}


\overline{A} z > 0,\quad \overline{B} z < 0,


\end{equation}


где $\overline{A}:= [a_{js}]$, $a_{js}:= f_j(a_s)$;


$\overline{B}:= [b_{is}]$, $b_{is}:= f_s(b_i)$.





Если $\Bar{z} = [\nk{\Bar{z_1}}{\Bar{z_k}}]$ ---


некоторое решение системы (7.2), то функция $f(x) =\svij{s=1}{k}


\Bar{z}_s\, f_s(x)$ будет строго разделяющей множества $A$ и


$B$. Дискриминантная функция $f(x)$ реализует соотнесение


(по свойству принадлежности одному из образов ${\cal G}$ и


${\cal L}$) предъявляемого для распознавания вектора $y$ по правилу


$$


y \in {\cal G},\quad \text{если}\quad f(y) > 0;\qquad


y \in {\cal L},\quad \text{если}\quad f(y) < 0.


$$





Функция $f(x)$ реализует {\it решающее правило}.


Система (7.2) может быть как совместной, так и несовместной.


На случай несовместности имеется обобщение понятия решения как


конечной совокупности векторов $\{c_l\} \su\R^n$ (именуемой


{\it комитетным решением}) такой, что любое из неравенств


системы (7.2) удовлетворяется более чем половиной векторов этой


совокупности. Комитетная технология и ее использования в задачах


распознавания составляют важное и хорошо разработанное направление в


распознавании образов \cite{Maz}.





Если $M$ и $N$ --- многогранники, причем $M \,\bigcap\,


N = \emp$, то эти множества строго разделяются аффинной функцией.


Ниже речь пойдет о разделимости произвольных непересекающихся


полиэдральных множеств кусочно-линейной функцией ($k$-функцией).





Итак, пусть $M = \,\bigcup\limits_{(j)}\, M_j$, $N =


\,\bigcup\limits_{(i)}\, N_i$, $\{M_j\}$ и $\{N_i\}$ ---


конечные совокупности выпуклых многогранников, причем $M


\,\bigcap\, N = \emp$. Имеет место





\begin{thms}


Полиэдральные множества $M$ и $N$ с пустым пересечением


строго разделяются $k$-функцией, т.е. функцией вида {\em (2.2)}


\begin{equation}


f (x) = \min_{j \in \overline{1,m}}\, |A_j x -b^j|_{\max}.


\end{equation}


\end{thms}





\begin{pf} Полиэдральное множество может быть задано одним


$k$-неравенством: если $M_j = \{x \mid A_j x \le b^j \}$ и


$N_i = \{ x \mid B_i x \ge d^i \}$, то


$$


M = \{ x \mid f (x) \le 0 \},


$$


где $f(x)$ --- функция (7.3);


$$


N = \{ x \mid g (x) \ge 0 \},


$$


где $g (x) =\max\limits_{(i)} |B_i x -d^i|_{\min}$.


Учитывая соотношения: $\bold X\backslash M \supset N$,


$\bold X\backslash N \supset M$, при любых $\al > 0$ и


$\be > 0$ будем иметь


$$


\al\, f(x) + \be\, g(x) < 0\; \fo\, x \in M;\quad


\al\, f(y) + \be\, g(y) > 0\; \fo\, y \in N.


$$


Таким образом, построена функция $f_{\alpha,\beta} (x):=


\al\, f(x) + \be\, g(x)$, зависящая от числовых параметров $\al >


0$ и $\be > 0$, строго разделяющая полиэдральные множества $M$


и $N$.


\end{pf}





\begin{cors}


Пусть $\R^n \supset \{a_j\}_{j \in J}$ и


$\R^n \supset \{b_i\}_{i \in I}$ --- конечные совокупности точек


из $\R^n$, при этом $\{a_j\} \,\bigcap\,


\{b_i\} = \emp$. Тогда существует $k$-функция $f(x)$, строго


разделяющая эти множества, т.е.


$$


f(a_j) < 0\; \fo\, j \in J;\quad f(b_i) > 0\; \fo\, i \in I.


$$


\end{cors}





Действительно,


поскольку точка из $\R^n$ может быть задана конечной системой


линейных неравенств (если $\Bar x = [


\nk{\Bar{x}_1}{\Bar{x}_n}]$, то $\Bar x = \{ x \mid


\Bar x \le x \le \Bar x \}$), то сформулированное следствие


превращается в частный случай теоремы 7.1.


На самом деле, в данном следствии пространство, из которого


берутся точки, может быть произвольным. Сведение к


конечномерному арифметическому


пространству тривиально: достаточно взять


подпространство $\En{k}$ исходного пространства, натянутое на


совокупность $\{a_j\}_{j \in J} \,\bigcup\, \{b_i\}_{i \in I}$,


выделить его базис, в котором векторы $a_j$ и $b_i$


будут представлены точками конечномерного


арифметического пространства. Если взять прямое


разложение $\bold X=\En{k}+\bold H$, то $k$-функция $f(x)$,


разделяющая указанные совокупности точек в $\En{k}$, может быть


продолжена до $k$-функции $\Bar{f}(x)$, определенной на всем


пространстве $\bold X$, по правилу: если $x \in\bold X$ и $x


= \Bar{x} + h$, $\Bar{x} \in\En{k}$, $h \in\bold H$,


то $\Bar{f}(x) =f(\Bar{x})$.








\begin{thebibliography}{9} 





\bibitem{Plo}


Плотников С.В. {\em Методы проектирования в задачах нелинейного


программирования}. Дисс.$\;\dotsc$ канд. физ.-матем. наук. -- Свердловск:


УрГУ, 1983. -- С.\,83--118.


\bibitem{Vol}


Волокитин Е.П. {\em О представлении непрерывных кусочно-линейных


функций} // Управляемые системы. -- Новосибирск: Наука, 1979.


-- \No\,19. -- C.\,14--21.


\bibitem{Mel}


Meltzer D. {\em On the expressibility of piecewise linear


continuous functions as the difference of two piecewise linear


convex functions} // Math. Program., Study 29. -- 1986.


-- P.\,118--134.


\bibitem{Kri}


Kripfganz A., Schulze R. {\em Piecewise affine functions as


a difference of two convex functions} // Optimization. -- 1987.


-- V.\,18. -- \No\,1. -- P.\,23--29.


\bibitem{Ben}


Benchekroun B. {\em A nonconvex piecewise linear optimization


problem} // Computers Math. Applic. -- 1991.


-- V.\,21. -- \No\,6\,/\,7. -- P.\,77--85.


\bibitem{Gor}


Gorokhovik V.V., Zorko O.I. {\em Piecewise affine functions


and polyhedral sets} // Optimization. -- 1994.


-- V.\,33. -- P.\,209--221.


\bibitem{AHU}


{\em Исследования по линейному и нелинейному программированию} /


Под ред. K.Дж.\,Эрроу, Д.\,Гурвица, Х.\,Удзавы.


-- M.: Иност. Лит., 1962. -- 298~c.


\bibitem{Maz}


Мазуров В.Д. {\em Метод комитетов в задачах оптимизации и


классификации}. -- M.: Наука, 1990. -- 246~c.





\end{thebibliography} 


 


\vskip 2cm 


 


\post{\it Институт математики и механики}{\it Поступила}


\post{\it Уральского отделения}{22.08.1997} 


\post{\it Российской Академии Наук}{} 





\end{document}








